% !TeX root=../main.tex
\chapter{مبانی نظری و پیشینه پژوهش}
\label{chap:background}

\section{مقدمه}
در این فصل ابتدا مبانی علمی لازم برای درک کدهای تصحیح خطا و به‌ویژه کدهای توربو مرور می‌شود؛ سپس مفهوم شناسایی کور و پژوهش‌های مرتبط در کدهای کانولوشنی و توربو دسته‌بندی می‌گردد و در پایان شکاف پژوهشی مشخص می‌شود.

\section{کدهای تصحیح خطا}
\subsection{مفهوم کدگذاری کانال}
کدگذاری کانال با افزودن افزونگی کنترل‌شده، امکان تشخیص و تصحیح خطا در گیرنده را فراهم می‌کند.

\paragraph{کدهای بلوکی}
در کدهای بلوکی، هر بلوک اطلاعات مستقل رمزگذاری می‌شود و ساختار جبری مشخصی دارد.

\paragraph{کدهای کانولوشنی}
کدهای کانولوشنی با حافظهٔ محدود و نمایش ترلیس، پایهٔ بسیاری از کدگذارهای عملی از جمله RSC هستند.

\paragraph{کدهای LDPC}
کدهای بررسی توازن کم‌چگال با گراف دوقسمتی و رمزگشایی تکراری شناخته می‌شوند.

\paragraph{کدهای توربو}
کدهای توربو با ترکیب موازی یا سریال کدگذارهای کانولوشنی و درهم‌ریز، عملکرد نزدیک به ظرفیت شانون را ممکن می‌سازند.

\section{کدهای توربو}
\subsection{تاریخچه و تکامل کدهای توربو}
پیدایش کدهای توربو نقطهٔ عطفی در نظریهٔ کدگذاری بود
\cite{Berrou93}
و به استانداردهای مخابراتی متعددی راه یافت.

\subsection{ساختار Turbo Code}
ساختار متداول شامل دو کدگذار RSC، درهم‌ریز داخلی، و در برخی سامانه‌ها درهم‌ریز خارجی و پانچینگ است.

\subsection{کدگذارهای RSC}
کدگذارهای سیستماتیک بازگشتی پایهٔ بسیاری از طرح‌های توربو هستند و رفتار ترلیس و قیود پاریتی آن‌ها در بازسازی کور نقش کلیدی دارد.

\subsection{پلی‌نوم‌های مولد}
پلی‌نوم‌های بازخورد و پیش‌خور ساختار کدگذار را تعیین می‌کنند و یکی از اهداف اصلی بازسازی کور هستند.

\subsection{حافظه و حالت‌های کدگذار}
حافظهٔ کدگذار تعداد بیت‌های رجیستر و تعداد حالت‌های ترلیس را مشخص می‌کند.

\subsection{ترلیس و گذارهای حالت}
نمایش ترلیس گذارهای مجاز حالت را توصیف می‌کند و مبنای رمزگشایی و استخراج قیود است.

\subsection{خاتمه ترلیس}
بیت‌های خاتمه، ترلیس را به حالت مشخصی می‌برند و طول فریم خروجی را تحت تأثیر قرار می‌دهند؛ در مدل این رساله رابطهٔ طول با خاتمه لحاظ می‌شود.

\subsection{درهم‌ریز داخلی}
درهم‌ریز داخلی ترتیب بیت‌های ورودی کدگذار دوم را تغییر می‌دهد و بازیابی آن از مسائل دشوار بازسازی است.

\subsection{درهم‌ریز خارجی}
درهم‌ریز خارجی بر چینش سریالی خروجی اثر می‌گذارد و باید همراه با سایر پارامترها بازسازی شود.

\subsection{پانچینگ و نرخ کد}
پانچینگ با حذف گزینشی بیت‌های پاریتی نرخ کد را افزایش می‌دهد و شناسایی آن را دشوارتر می‌کند.

\subsection{ساختار سریالی خروجی کد توربو}
چینش بیت‌های سیستماتیک و پاریتی در جریان خروجی، مدل مشاهدهٔ گیرنده را شکل می‌دهد.

\section{مدل کانال و مدل مشاهده}
\subsection{کانال AWGN}
مدل نویز سفید گاوسی جمع‌شونده فرض پایهٔ بسیاری از ارزیابی‌ها است.

\subsection{BPSK}
مدولاسیون دودویی فاز، نگاشت بیت به نماد را ساده می‌کند و با مدل‌های تحلیلی سازگار است.

\subsection{آشکارسازی سخت}
در آشکارسازی سخت، هر نماد به یک بیت قطعی نگاشت می‌شود و مشاهده به دنبالهٔ باینری تبدیل می‌گردد.

\subsection{آشکارسازی نرم}
اطلاعات نرم (مانند LLR) می‌تواند دقت بازسازی را بهبود دهد و در فصل توسعه بررسی می‌شود.

\subsection{مدل کانال معادل BSC}
پس از آشکارسازی سخت، کانال معادل BSC با احتمال خطای بیت مشخص، مدل مشاهدهٔ عملیاتی رساله را تشکیل می‌دهد.

\section{مفهوم شناسایی کور}
در ادبیات، تمایز میان طبقه‌بندی کور، تخمین کور و \gls{blind-recon} اهمیت دارد. طبقه‌بندی نوع کد را مشخص می‌کند، تخمین برخی پارامترها را به‌دست می‌آورد، و بازسازی ساختار کامل (یا مجموعهٔ ساختارهای مجاز) را هدف قرار می‌دهد؛ تمرکز این رساله بر بازسازی است. کدهای \gls{turbo-code} و کدگذارهای \lr{RSC} نمونهٔ مرکزی این مطالعه هستند.

\section{شناسایی کور کدهای کانولوشنی}
روش‌های مبتنی بر رتبه، قیود دوگان، و الگوریتم‌های تکراری برای بازیابی پارامترهای کد کانولوشنی مرور می‌شوند تا زمینه‌ای برای تعمیم به توربو فراهم گردد.

\section{شناسایی کور کدهای توربو}
\subsection{شناسایی پارامترهای کدگذار}
\subsection{بازیابی پلی‌نوم‌ها}
\subsection{بازیابی درهم‌ریز}
\subsection{تشخیص نرخ کد}
\subsection{تشخیص طول فریم}
\subsection{تشخیص حافظه}
\subsection{تشخیص پانچینگ}
\subsection{همزمان‌سازی و تعیین مرز فریم}
در این زیربخش‌ها پژوهش‌های مرتبط با هر پارامتر دسته‌بندی و نقاط قوت و ضعف آن‌ها بیان می‌شود.

\section{روش‌های موجود در ادبیات}
به‌جای فهرست خطی مقالات، روش‌ها در دسته‌هایی مانند شناسایی کدگذار، بازیابی درهم‌ریز، dual-code، BCJR، rank-based، ML/EM و joint reconstruction طبقه‌بندی می‌شوند. جدول~\ref{tab:lit-categories} نمای کلی این دسته‌بندی را نشان می‌دهد.

\begin{table}[ht]
	\caption{دسته‌بندی روش‌های ادبیات در شناسایی/بازسازی کور کد توربو}
	\label{tab:lit-categories}
	\centering
	\begin{tabularx}{0.95\textwidth}{|X|X|}
		\hline
		\textbf{دسته} & \textbf{پارامتر بازیابی‌شدهٔ متداول} \\
		\hline
		شناسایی کدگذار & پلی‌نوم‌ها، حافظه \\
		\hline
		بازیابی درهم‌ریز & جایگزشت \\
		\hline
		dual-code & قیود پاریتی \\
		\hline
		مبتنی بر BCJR & ساختار و درهم‌ریز \\
		\hline
		rank-based & طول، بعد، ساختار \\
		\hline
		ML/EM & تخمین پارامتر \\
		\hline
		بازسازی مشترک & چند مؤلفهٔ همزمان \\
		\hline
	\end{tabularx}
\end{table}

\section{شکاف پژوهشی}
با وجود پیشرفت‌های قابل توجه، بازسازی همزمان ساختار کامل یک کد توربو شامل پارامترهای کدگذار، درهم‌ریز داخلی و درهم‌ریز خارجی، در شرایط مشاهدهٔ ناقص و نویزی، همچنان مسئله‌ای دشوار و تا حد زیادی حل‌نشده است. این رساله همین خلأ را هدف قرار می‌دهد.
